\chapter{答案-几何}

\section{余角补角}

\subsection{与物理结合}
\label{几何-场景-题目1}

由题意，入射角与折射角之比为 \(4:3\)，即
\[
\frac{\angle\text{入射}}{\angle\text{折射}} = \frac{4}{3} \quad\Longrightarrow\quad \angle\text{折射} = \frac{3}{4}\,\angle\text{入射}.
\]

设两条入射光线的入射角分别为 \(\alpha\) 和 \(\beta\)，则它们对应的折射角分别为
\[
\theta_1 = \frac{3}{4}\alpha,\qquad \theta_2 = \frac{3}{4}\beta.
\]

观察题图可知，两条入射光线位于法线两侧（即分别从法线左右射入），因此折射光线也分别位于法线两侧，两折射光线的夹角 \(\gamma\) 等于两个折射角之和：
\[
\gamma = \theta_1 + \theta_2 = \frac{3}{4}\alpha + \frac{3}{4}\beta = \frac{3}{4}(\alpha+\beta).
\]

所以正确的关系为 \(\dfrac{3}{4}(\alpha+\beta)=\gamma\)，对应选项 \textbf{B}。

\vspace{2mm}
\noindent 答案：\textbf{B}


\subsection{几何-比例-题目1}
\label{几何-比例-题目1}

\begin{center}
\includegraphics[width=1.0\textwidth]{几何-比例-题目1-答案.png}
\end{center}


\subsection{几何-比例-题目2}
\label{几何-比例-题目2}

\begin{center}
\includegraphics[width=1.0\textwidth]{几何-比例-题目2-答案.png}
\end{center}